\section{The induced completion and when it is proper}
\label{sec:completion}

The finite-complement extensions of the preceding section define a descending
family of equivalence relations.  Write $x\stageeq{n}y$ when no compatible
observer using at most $n+1$ states outside the base distinguishes $x$ from
$y$.  Increasing the budget can only refine the relation, and
\cref{thm:finite-external-separation} implies
\[
  \bigcap_{n\ge 0}\stageeq{n}\;=\;=.
\]
Thus the resulting filtered space is separated.

A sequence $(x_k)$ is called Cauchy when, for every fixed $n$, all sufficiently
late terms are pairwise $\stageeq{n}$-equivalent.  Two Cauchy sequences are
identified when they eventually agree at every fixed stage.  This gives the
usual completion of a separated filtration; constant sequences embed the
original generated algebra injectively.

The question of interest is whether this embedding is surjective.  We call the
completion \emph{proper} when it contains a Cauchy class not represented by a
single generated term.

\subsection{A finite-state lemma}

Consider a sequence obtained by repeatedly applying one unary context.  At a
fixed observation stage, its interpretation is an orbit
\[
  s_{k+1}=F(s_k)
\]
of a deterministic self-map.  Suppose the part of this orbit relevant at that
stage can be encoded by $N$ states, in the sense that equality of codes at two
ordered times forces equality of the corresponding interpreted states.

\begin{lemma}[Factorial synchronization]
\label{lem:factorial-synchronization}
Under the hypothesis above, the factorial subsequence
\[
  s_{0!},s_{1!},s_{2!},\ldots
\]
is eventually constant at that observation stage.
\end{lemma}

\begin{proof}
Among the first $N+1$ coded orbit positions, two have the same code.  Hence the
orbit is eventually periodic with a period $p\le N$.  Every sufficiently large
factorial is divisible by $p$, so all sufficiently late factorial indices lie
in the same position of that eventual cycle.
\end{proof}

This is the only finite-state input used below.

\subsection{A growing orbit cannot converge to a finite term}

Assume now that
\[
  t_{k+1}=\susp(u,t_k,\old(e))
\]
for a fixed operation symbol $u$ and fixed base element $e$, and that every
step is genuinely unresolved, so the constructor size of $t_k$ strictly
increases.

Fix any generated term $x$.  Choose $N$ so large that $t_{N+1}$ is larger than
every subterm of $x$ and $t_N$.  Apply the finite-family construction of
\cref{thm:finite-pattern-realization} to the two roots $x$ and $t_N$.  The
observer remembers every selected subterm exactly.  The transition producing
$t_{N+1}$ is not among those selected transitions, so it is sent to the sink
state $\overflow$.  By construction the same unary context keeps
$\overflow$ at $\overflow$.  Therefore all later $t_k$ have the same sink
value, while $x$ retains its distinct selected value.

\begin{lemma}[No generated limit]
\label{lem:no-generated-limit}
For every generated term $x$, there is a finite-complement observer and an
index $N$ such that the observer separates $x$ from every $t_k$ with $k>N$.
Consequently no subsequence of $(t_k)$ that is eventually contained in this
tail can converge to $x$ in the filtered completion.
\end{lemma}

The argument is elementary, but it is important that the observer preserves the
entire base pointwise while using only finitely many new states.

\subsection{A sufficient criterion for properness}

Combining the two preceding lemmas yields the general mechanism used in the
examples.

\begin{theorem}[Properness from finite orbit compression]
\label{thm:compression-escape-properness}
Let $(t_k)$ be a growing orbit of the form above.  Suppose that, at every
observation stage, its interpreted trajectory admits a finite code.  Then the
factorial subsequence $(t_{k!})$ is Cauchy and has no generated limit.  Hence
the completion is proper.
\end{theorem}

\begin{proof}
At each fixed stage, \cref{lem:factorial-synchronization} makes the factorial
subsequence eventually constant, so it is Cauchy.  By
\cref{lem:no-generated-limit}, every generated candidate is separated from the
whole sufficiently late tail by some observer.  Thus no generated term
represents its Cauchy class.
\end{proof}

\subsection{Finite bases}

Suppose that $A$ injects into $\Fin(m)$ and that some application
$\eval_D(f,a,z)$ is undefined.  Starting with $c_0=\old(a)$, define
\[
  c_{k+1}=f(c_k,\old(z)).
\]
At stage $n$, the whole observer has at most $m+n+1$ states after coding the
base, so the interpreted orbit is finitely coded.  The syntax grows by one
unresolved constructor at each step.  The preceding theorem therefore gives a
new completion point.

\begin{theorem}[Finite-base criterion]
\label{prop:proper-completion}
If the base is finitely coded, the completion is proper if and only if the
original evaluator is not total.
\end{theorem}

\begin{proof}
If some base application is undefined, the construction above gives a Cauchy
sequence with no generated limit.  Conversely, if the evaluator is total, every
generated expression reduces to a base element.  Since the filtration separates
base elements, every Cauchy sequence is eventually constant and no new point is
added.
\end{proof}

\subsection{An infinite-base criterion}

The preceding argument does not require the whole base to be finite if the
chosen orbit cannot use base elements as additional memory.  Assume
\[
  \eval_D(g,a,b)=\operatorname{none}
\]
and suppose there are an operation $u$ and a base element $e$ such that
\[
  \eval_D(u,x,e)=\operatorname{some}(x)
  \qquad (x\in A).
\]
Let
\[
  c_0=\susp(g,\old(a),\old(b)),
  \qquad
  c_{k+1}=\susp(u,c_k,\old(e)).
\]
In a stage-$n$ observer, if the interpreted orbit ever enters the base, the
identity law above makes it constant from then on.  Otherwise it remains among
the $n$ tags and the overflow state.  Hence the trajectory always has a finite
code of size at most $n+2$.

\begin{theorem}[Base-fixing criterion]
\label{thm:old-fixing-proper-completion}
Under the preceding hypotheses, the completion is proper.  Moreover, for every
$n$ there are distinct iterates that cannot be separated at stage $n$; in
particular the least separation budgets along the orbit are unbounded.
\end{theorem}

\begin{proof}
The properness statement follows from
\cref{thm:compression-escape-properness}.  For the quantitative statement, the
finite code has at most $n+1$ non-base states unless the orbit has already
stabilized in the base.  Factorial synchronization therefore gives
\[
  c_{(n+1)!}\stageeq{n}c_{(n+2)!}.
\]
The two terms are syntactically distinct because each iteration adds one
constructor.
\end{proof}

The finite-base and base-fixing hypotheses are independent.  A finite partial
algebra need not have a unary context fixing every base element, while natural
arithmetic has the context $x\mapsto x+0$ but has an infinite base.

\subsection{Comparison with the prefix-depth metric}

The filtration above is not merely the usual prefix metric on trees.  Let
\[
  b_0=\old(a),\qquad b_{k+1}=\susp(f,b_k,\old(z)),
\]
with $\eval_D(f,a,z)$ undefined.  For every fixed syntactic depth, the sequence
$(b_k)$ is eventually constant on that prefix and is therefore Cauchy in the
standard tree metric.

By contrast, a single one-tag finite-complement observer can be chosen so that
the comb context alternates between its tag and the overflow state.  The same
observer then distinguishes every adjacent pair $b_k,b_{k+1}$, regardless of
how deep the new constructor occurs.

\begin{proposition}[The two Cauchy notions differ]
\label{prop:depth-observation-contrast}
There exists a generated comb sequence that is Cauchy for prefix-depth
agreement but is not Cauchy for the finite-complement filtration.
\end{proposition}

Thus the completion studied here is not simply the ordinary infinite-tree
completion written in different notation.
